\chapter{The planted-forest residue \texorpdfstring{$L_\infty$}{L-infinity}-algebra of \texorpdfstring{$\FM_3(X)$}{FM3(X)}}
\label{app:fm3-planted-forest-residue}

\index{Fulton--MacPherson compactification!$\FM_3(X)$|textbf}
\index{planted-forest residue $L_\infty$-algebra|textbf}
\index{Maurer--Cartan obstruction!planted-forest cubic obstruction|textbf}

This appendix isolates the smallest Fulton--MacPherson compactification on which
all three layers of the modular mechanism are simultaneously visible:
\begin{enumerate}[label=\textup{(\roman*)}]
\item the local logarithmic counterterm ideal is contractible;
\item the surviving quotient is residue-theoretic; and
\item the first genuinely higher bracket already records a nested collision.
\end{enumerate}
The global cyclic deformation package is already established in
Theorems~\ref{thm:cyclic-linf-graph}, \ref{thm:universal-theta},
and~\ref{thm:mc2-full-resolution}.  The point here is different:
one wants a small, explicit, stratum-level model in which the nested-collision
mechanism can be written down on the nose.  For $\FM_3(X)$ one can do this
completely.

Everything in this appendix is at the \emph{model level} in the sense of the
H/M/S discipline: the objects are explicit dg and $L_\infty$ models built from
logarithmic differential forms on boundary strata and from the planted-forest
incidence relations of $\FM_3(X)$.

\section{Geometry of \texorpdfstring{$\FM_3(X)$}{FM3(X)}}
\label{sec:fm3-geometry-residue}

Let $X$ be a smooth complex curve.  Write
\[
\Delta_{ij} := \{x_i = x_j\} \subset X^3,
\qquad
\Delta_{123} := \{x_1 = x_2 = x_3\}.
\]
The Fulton--MacPherson compactification $\FM_3(X)$ is obtained by blowing up the
small diagonal and then the proper transforms of the pair diagonals.

\begin{proposition}[Boundary geometry of \texorpdfstring{$\FM_3(X)$}{FM3(X)}; \ClaimStatusProvedHere]
\label{prop:fm3-boundary-geometry}
Let $X$ be a smooth complex curve.  Then:
\begin{enumerate}[label=\textup{(\alph*)}]
\item $\FM_3(X)$ is smooth and its boundary is a simple normal crossing divisor
  with irreducible components
  \[
  D_{12},\; D_{23},\; D_{13},\; D_{123},
  \]
  where $D_{ij}$ denotes the proper transform of $\Delta_{ij}$ and
  $D_{123}$ denotes the exceptional divisor over $\Delta_{123}$.
\item The divisors $D_{12},D_{23},D_{13}$ are pairwise disjoint.
\item Each $D_{ij}$ meets $D_{123}$ transversely along a smooth codimension-two
  stratum
  \[
  S_{ij} := D_{ij}\cap D_{123}.
  \]
\item Every boundary stratum is one of
  \[
  D_{12},\ D_{23},\ D_{13},\ D_{123},\ S_{12},\ S_{23},\ S_{13}.
  \]
\item One has canonical identifications
  \[
  D_{ij} \cong X^2,
  \qquad
  D_{123} \cong \mathbf{P}(N_{\Delta_{123}/X^3}),
  \qquad
  S_{ij} \cong X,
  \]
  and each $S_{ij}\to X$ is a section of the $\mathbf{P}^1$-bundle
  $D_{123}\to X$.
\end{enumerate}
\end{proposition}

\begin{proof}
The statement is local on $X$.  Choose a coordinate $z$ on a disk in $X$ and
use coordinates
\[
(x,u,v) := (z_1, z_2-z_1, z_3-z_1)
\]
on $X^3$ near the small diagonal.  Then
\[
\Delta_{123} = \{u=v=0\},
\qquad
\Delta_{12} = \{u=0\},
\qquad
\Delta_{13} = \{v=0\},
\qquad
\Delta_{23} = \{u-v=0\}.
\]
Blowing up $\Delta_{123}$ amounts in the normal slice to blowing up the origin
in $\mathbf{A}^2_{u,v}$.  The exceptional divisor is the projectivized normal
space, hence a copy of $\mathbf{P}^1$ over each point of the small diagonal.
The strict transforms of the three lines $u=0$, $v=0$, and $u-v=0$ in
$\operatorname{Bl}_0 \mathbf{A}^2$ are pairwise disjoint and meet the
exceptional divisor at the three distinct points
$[0:1],\ [1:0],\ [1:1]$
in the exceptional $\mathbf{P}^1$.  Therefore, after the second-stage blowups
along the strict transforms of the pair diagonals, the resulting total boundary
is simple normal crossings, the three divisors $D_{12},D_{23},D_{13}$ are
pairwise disjoint, and each meets the exceptional divisor $D_{123}$
transversely along a section.

The description of the strata follows immediately: there are four codimension-one
strata and three codimension-two strata, with no higher intersections.  Since a
pair diagonal on a curve has codimension one, its projectivized normal direction
is a point, so the boundary divisor $D_{ij}$ is canonically isomorphic to the
base $\Delta_{ij} \cong X^2$.  The exceptional divisor over the small diagonal is
$\mathbf{P}(N_{\Delta_{123}/X^3})$, a $\mathbf{P}^1$-bundle over $X$.  Each
intersection $S_{ij}$ is one of its three distinguished sections.  Globalizing
from the local chart proves the claim.
\end{proof}

\begin{definition}[The planted forests of \texorpdfstring{$\FM_3(X)$}{FM3(X)}]
\label{def:fm3-planted-forests}
The boundary stratification of $\FM_3(X)$ is indexed by the seven planted forests
\[
12,\ 23,\ 13,\ 123,\ 12\prec 123,\ 23\prec 123,\ 13\prec 123.
\]
Here $ij$ denotes the binary collision of points $i$ and $j$, $123$ denotes the
triple collision, and $ij\prec 123$ denotes the nested collision in which the
binary cluster $ij$ sits inside the triple cluster $123$.
\end{definition}

\begin{remark}[Why $\FM_3$ is the first nontrivial atom]
On $\FM_2(X)$ there is only one boundary divisor and no nested boundary.
Consequently the residue quotient is a strict dg Lie algebra and no higher
transferred operation can appear.  The space $\FM_3(X)$ is the first case in
which a codimension-two stratum survives after the local counterterm ideal is
quotiented out; exactly there the first higher operation $l_3$ appears.
\end{remark}

\section{The boundary residue dg Lie algebras}
\label{sec:fm3-residue-dglas}

Fix a complete filtered dg Lie algebra $\mathfrak{g}$.  In applications one may
specialize to $\mathfrak{g}=\Defcyc(\cA)$ or to one of its explicit dg models.
Choose a $C^\infty$ logarithmic de~Rham model on $\FM_3(X)$; this is enough for
the model-level constructions of the present appendix.

At the generic point of $D_{ij}$ choose a local defining equation $r_{ij}=0$ and
write
\[
\eta_{ij} := d\log r_{ij}.
\]
At the generic point of $D_{123}$ choose a local defining equation
$r_{123}=0$ and write
\[
\eta_{123} := d\log r_{123}.
\]
Near the codimension-two stratum $S_{ij}$ one may choose local coordinates
$(r_{ij},r_{123},y)$ so that
\[
D_{ij} = \{r_{ij}=0\},
\qquad
D_{123} = \{r_{123}=0\},
\qquad
S_{ij} = \{r_{ij}=r_{123}=0\},
\]
and the logarithmic forms are generated by $\eta_{ij}$ and $\eta_{123}$.

Define the residue dg Lie algebras
\begin{align*}
\mathfrak{r}_{12} &:= \Omega^\bullet(D_{12})\,\widehat\otimes\,\Lambda\langle \eta_{12}\rangle\,\widehat\otimes\,\mathfrak{g},\\
\mathfrak{r}_{23} &:= \Omega^\bullet(D_{23})\,\widehat\otimes\,\Lambda\langle \eta_{23}\rangle\,\widehat\otimes\,\mathfrak{g},\\
\mathfrak{r}_{13} &:= \Omega^\bullet(D_{13})\,\widehat\otimes\,\Lambda\langle \eta_{13}\rangle\,\widehat\otimes\,\mathfrak{g},\\
\mathfrak{r}_{123} &:= \Omega^\bullet(D_{123})\,\widehat\otimes\,\Lambda\langle \eta_{123}\rangle\,\widehat\otimes\,\mathfrak{g},
\end{align*}
and on the codimension-two strata
\begin{align*}
\mathfrak{s}_{12} &:= \Omega^\bullet(S_{12})\,\widehat\otimes\,\Lambda\langle \eta_{12},\eta_{123}\rangle\,\widehat\otimes\,\mathfrak{g},\\
\mathfrak{s}_{23} &:= \Omega^\bullet(S_{23})\,\widehat\otimes\,\Lambda\langle \eta_{23},\eta_{123}\rangle\,\widehat\otimes\,\mathfrak{g},\\
\mathfrak{s}_{13} &:= \Omega^\bullet(S_{13})\,\widehat\otimes\,\Lambda\langle \eta_{13},\eta_{123}\rangle\,\widehat\otimes\,\mathfrak{g}.
\end{align*}
Each of these is a complete filtered dg Lie algebra with differential
$d_{\mathrm{dR}}\otimes 1 + 1\otimes d_{\mathfrak g}$ and bracket given by wedge
product on forms and the Lie bracket on $\mathfrak g$.

Set
\[
\mathfrak{r}_0 := \mathfrak{r}_{12}\oplus\mathfrak{r}_{23}\oplus\mathfrak{r}_{13}\oplus\mathfrak{r}_{123},
\qquad
\mathfrak{r}_1 := \mathfrak{s}_{12}\oplus\mathfrak{s}_{23}\oplus\mathfrak{s}_{13}.
\]
There is a natural dg Lie morphism
\[
\chi \colon \mathfrak{r}_0 \longrightarrow \mathfrak{r}_1
\]
given by the planted-forest incidence relations:
\begin{equation}
\label{eq:chi-fm3}
\chi(\alpha_{12},\alpha_{23},\alpha_{13},\alpha_{123})
:=
\bigl(
\alpha_{12}|_{S_{12}}-\alpha_{123}|_{S_{12}},
\alpha_{23}|_{S_{23}}-\alpha_{123}|_{S_{23}},
\alpha_{13}|_{S_{13}}-\alpha_{123}|_{S_{13}}
\bigr).
\end{equation}

\begin{lemma}[The planted-forest incidence map; \ClaimStatusProvedHere]
\label{lem:chi-dgla-fm3}
The map $\chi$ of~\eqref{eq:chi-fm3} is a filtered dg Lie morphism.
\end{lemma}

\begin{proof}
Restriction of differential forms to a closed stratum commutes with the de~Rham
differential and preserves wedge products.  The coefficient bracket on
$\mathfrak g$ is untouched by restriction.  Therefore each component of $\chi$
commutes with the differential and preserves brackets; completeness and the
filtration are inherited componentwise.
\end{proof}

\section{The small planted-forest model}
\label{sec:fm3-small-model}

Let $U_{12},U_{23},U_{13},U_{123}$ be disjoint tubular neighborhoods of the four
boundary divisors chosen so that the only nonempty overlaps are
\[
U_{12,123}:=U_{12}\cap U_{123},
\qquad
U_{23,123}:=U_{23}\cap U_{123},
\qquad
U_{13,123}:=U_{13}\cap U_{123},
\]
with no triple overlaps.  Let $\mathfrak{L}_\partial(\mathfrak g)$ denote the
logarithmic dg Lie algebra on the resulting boundary neighborhood.
By the local contraction theorem proved earlier in the thread, the positive
normal-weight ideal on each patch and overlap is contractible, so the patchwise
logarithmic dg Lie algebra retracts to the corresponding residue dg Lie algebra.
Because the \v{C}ech nerve of this cover is truncated at degree~$1$, the
residue semicosimplicial dg Lie algebra is precisely the dg Lie morphism
$\chi\colon \mathfrak r_0\to \mathfrak r_1$.

\begin{definition}[The planted-forest residue complex]
\label{def:fm3-planted-forest-complex}
Define
\[
\mathfrak{P}_{\FM_3}(X,\mathfrak g) := \operatorname{Cone}(\chi)[-1].
\]
Thus, in degree $k$,
\[
\mathfrak{P}_{\FM_3}(X,\mathfrak g)^k
=
\mathfrak r_0^k \oplus \mathfrak r_1^{k-1}.
\]
We write an element as $(\alpha,\beta)$ with
$\alpha\in \mathfrak r_0$ and $\beta\in \mathfrak r_1[-1]$.
The cone differential is
\begin{equation}
\label{eq:l1-fm3-cone}
 l_1(\alpha,\beta) := \bigl(d\alpha,\, \chi(\alpha)-d\beta\bigr).
\end{equation}
\end{definition}

\begin{theorem}[Boundary logarithmic dg Lie algebra as a homotopy fiber; \ClaimStatusProvedHere]
\label{thm:fm3-homotopy-fiber}
There is a filtered $L_\infty$ quasi-isomorphism
\[
\mathfrak{P}_{\FM_3}(X,\mathfrak g)
\rightsquigarrow
\mathfrak{L}_\partial(\mathfrak g).
\]
Equivalently, $\mathfrak{P}_{\FM_3}(X,\mathfrak g)$ is a small complete filtered
model for the boundary logarithmic deformation algebra of $\FM_3(X)$.
\end{theorem}

\begin{proof}
The chosen tubular cover has no nonempty triple overlaps, so its residue
semicosimplicial dg Lie algebra is the two-term diagram
\[
\mathfrak r_0 \xrightarrow{\chi} \mathfrak r_1.
\]
The Thom--Whitney totalization of such a semicosimplicial dg Lie algebra is the
standard path-space model of the homotopy fiber of $\chi$.  The normalized
complex of this path-space model is canonically the shifted cone
$\operatorname{Cone}(\chi)[-1]$.

Patchwise, the logarithmic dg Lie algebra splits as residue part plus positive
normal-weight ideal.  The positive normal-weight ideal is contractible on each
patch and on each overlap, hence the logarithmic semicosimplicial dg Lie algebra
is levelwise quasi-isomorphic to the residue semicosimplicial dg Lie algebra.
Totalization preserves filtered quasi-isomorphisms.  Finally,
Theorem~\ref{thm:htt} transfers the dg Lie structure of the Thom--Whitney
model to an $L_\infty$ structure on the shifted cone, yielding the desired
filtered $L_\infty$ quasi-isomorphism.
\end{proof}

\section{The explicit planted-forest \texorpdfstring{$L_\infty$}{L-infinity} brackets}
\label{sec:fm3-explicit-linf}

The shifted cone carries a canonical complete filtered $L_\infty$ structure.
Because the covering nerve has only degrees $0$ and $1$, the higher brackets are
as small as they can be: only one divisor input together with several overlap
inputs can produce a nonzero higher operation.

\begin{theorem}[Explicit planted-forest brackets for \texorpdfstring{$\FM_3(X)$}{FM3(X)}; \ClaimStatusProvedHere]
\label{thm:fm3-explicit-brackets}
The transferred $L_\infty$ structure on $\mathfrak P_{\FM_3}(X,\mathfrak g)$ has
nonzero brackets of the following forms.
\begin{enumerate}[label=\textup{(\alph*)}]
\item The unary bracket is the cone differential
  \[
  l_1(\alpha,\beta)=\bigl(d\alpha,\chi(\alpha)-d\beta\bigr).
  \]
\item For homogeneous $x=(\alpha,\beta)$ and $x'=(\alpha',\beta')$ with total
  degrees $|x|$ and $|x'|$ in the shifted cone,
  \begin{equation}
  \label{eq:l2-fm3}
  l_2(x,x')
  =
  \left(
  [\alpha,\alpha'],
  \frac12\Bigl([\beta,\chi(\alpha')] - (-1)^{|x||x'|}[\beta',\chi(\alpha)]\Bigr)
  \right).
  \end{equation}
\item The first higher bracket is supported exactly on the nested planted forests.
  For homogeneous overlap inputs $m_i=(0,\beta_i)$ and a homogeneous divisor
  input $u=(\alpha,0)$,
  \begin{equation}
  \label{eq:l3-fm3}
  l_3(m_1,m_2,u)
  =
  \frac1{12}
  \left(
  0,
  [\beta_1,[\beta_2,\chi(\alpha)]]
  + (-1)^{|m_1||m_2|}[\beta_2,[\beta_1,\chi(\alpha)]]
  \right),
  \end{equation}
  and all other values are obtained from~\eqref{eq:l3-fm3} by graded
  skew-symmetry.  Any input pattern not consisting of two overlap entries and one
  divisor entry vanishes.
\item More generally, for $n\geq 4$, the only potentially nonzero $l_n$ are those
  with one divisor input and $n-1$ overlap inputs; they are determined by the
  Bernoulli-series expansion of the gauge action in
  Theorem~\ref{thm:fm3-mc-dictionary} below.
\end{enumerate}
\end{theorem}

\begin{proof}
Existence of the transferred $L_\infty$ structure follows from
Theorem~\ref{thm:fm3-homotopy-fiber} together with the homotopy transfer theorem,
Theorem~\ref{thm:htt}.  The unary bracket is the cone differential.

To identify the first three operations explicitly, use the homotopy-fiber
interpretation.  Maurer--Cartan elements of the path-space model are pairs
consisting of a Maurer--Cartan element on level~$0$ and a gauge path in level~$1$
connecting $0$ to its image under~$\chi$.  The transferred Maurer--Cartan functor
is therefore the same as the Maurer--Cartan functor of the shifted cone.  The
resulting Maurer--Cartan equation is computed explicitly in
Theorem~\ref{thm:fm3-mc-dictionary}; expanding it to cubic order determines the
operations uniquely:
\begin{itemize}
\item the linear term gives $l_1$;
\item the quadratic terms are $\tfrac12[\alpha,\alpha]$ in the divisor sector and
  $\tfrac12[\beta,\chi(\alpha)]$ in the overlap sector, which is exactly
  \eqref{eq:l2-fm3};
\item the cubic term in the overlap sector is
  $\tfrac1{12}[\beta,[\beta,\chi(\alpha)]]$, forcing the coefficient
  $\tfrac1{12}$ in~\eqref{eq:l3-fm3}.
\end{itemize}
Because the cover has only degree-$1$ overlaps, there is no source for higher
operations with two divisor inputs or with overlap inputs supported on different
nested forests.  This proves the stated formulas.
\end{proof}

\begin{remark}[Componentwise form of \texorpdfstring{$l_3$}{l3}]
\label{rem:l3-componentwise-fm3}
The bracket $l_3$ is completely supported on the three codimension-two strata.
Writing
\[
\beta=(\beta_{12},\beta_{23},\beta_{13}) \in \mathfrak r_1,
\qquad
\alpha=(\alpha_{12},\alpha_{23},\alpha_{13},\alpha_{123}) \in \mathfrak r_0,
\]
one has
\begin{align*}
\bigl(l_3((0,\beta_1),(0,\beta_2),(\alpha,0))\bigr)_{12}
&=
\frac1{12}\Bigl([\beta_{1,12},[\beta_{2,12},\chi_{12}(\alpha)]]
+(-1)^{|m_1||m_2|}[\beta_{2,12},[\beta_{1,12},\chi_{12}(\alpha)]]\Bigr),\\
\bigl(l_3((0,\beta_1),(0,\beta_2),(\alpha,0))\bigr)_{23}
&=
\frac1{12}\Bigl([\beta_{1,23},[\beta_{2,23},\chi_{23}(\alpha)]]
+(-1)^{|m_1||m_2|}[\beta_{2,23},[\beta_{1,23},\chi_{23}(\alpha)]]\Bigr),\\
\bigl(l_3((0,\beta_1),(0,\beta_2),(\alpha,0))\bigr)_{13}
&=
\frac1{12}\Bigl([\beta_{1,13},[\beta_{2,13},\chi_{13}(\alpha)]]
+(-1)^{|m_1||m_2|}[\beta_{2,13},[\beta_{1,13},\chi_{13}(\alpha)]]\Bigr),
\end{align*}
where
\[
\chi_{12}(\alpha)=\alpha_{12}|_{S_{12}}-\alpha_{123}|_{S_{12}},
\quad
\chi_{23}(\alpha)=\alpha_{23}|_{S_{23}}-\alpha_{123}|_{S_{23}},
\quad
\chi_{13}(\alpha)=\alpha_{13}|_{S_{13}}-\alpha_{123}|_{S_{13}}.
\]
Thus $l_3$ is literally the nested planted-forest correction:
a binary residue $ij$ acts twice on the defect between the $ij$-residue and the
triple residue along the codimension-two stratum $ij\prec 123$.
\end{remark}

\section{Maurer--Cartan theory in closed form}
\label{sec:fm3-mc-theory}

The shifted-cone description is particularly effective because one can solve the
Maurer--Cartan equation in closed form.

\begin{theorem}[Maurer--Cartan dictionary for the planted-forest model; \ClaimStatusProvedHere]
\label{thm:fm3-mc-dictionary}
An element $(\alpha,\beta)\in \mathfrak r_0^1\oplus \mathfrak r_1^0$ is a
Maurer--Cartan element of $\mathfrak P_{\FM_3}(X,\mathfrak g)$ if and only if:
\begin{enumerate}[label=\textup{(\alph*)}]
\item $\alpha$ is a Maurer--Cartan element of $\mathfrak r_0$:
  \begin{equation}
  \label{eq:alpha-mc-fm3}
  d\alpha + \frac12[\alpha,\alpha] = 0;
  \end{equation}
\item the overlap parameter $\beta$ gauges $\chi(\alpha)$ to zero in
  $\mathfrak r_1$:
  \begin{equation}
  \label{eq:gauge-zero-fm3}
  e^{\beta} * \chi(\alpha) = 0.
  \end{equation}
\end{enumerate}
Equivalently,
\begin{equation}
\label{eq:beta-bernoulli-fm3}
 d\beta
 =
 \frac{\operatorname{ad}_{\beta}}{1-e^{-\operatorname{ad}_{\beta}}}\,\chi(\alpha)
 =
 \chi(\alpha)
 + \frac12[\beta,\chi(\alpha)]
 + \frac1{12}[\beta,[\beta,\chi(\alpha)]]
 - \frac1{720}[\beta,[\beta,[\beta,[\beta,\chi(\alpha)]]]]
 + \cdots.
\end{equation}
\end{theorem}

\begin{proof}
The path-space model of the homotopy fiber of $\chi$ identifies a Maurer--Cartan
element with the data of
\begin{enumerate}[label=\textup{(\roman*)}]
\item a Maurer--Cartan element $\alpha\in \mathfrak r_0$; and
\item a gauge path in $\mathfrak r_1$ connecting $0$ to $\chi(\alpha)$.
\end{enumerate}
Since the interval is contractible, every such path is gauge-equivalent to the
constant path determined by a single degree-zero parameter $\beta\in \mathfrak r_1^0$.
The endpoint condition is precisely~\eqref{eq:gauge-zero-fm3}.  This proves the
closed-form characterization.

To obtain the Bernoulli expansion, use the standard gauge action formula in a dg
Lie algebra:
\[
 e^{\beta} * x
 =
 e^{\operatorname{ad}_{\beta}}(x)
 - \frac{e^{\operatorname{ad}_{\beta}}-1}{\operatorname{ad}_{\beta}}(d\beta).
\]
Setting $x=\chi(\alpha)$ and imposing $e^{\beta}*x=0$ gives
\[
 d\beta
 =
 \frac{\operatorname{ad}_{\beta}}{e^{\operatorname{ad}_{\beta}}-1}
 e^{\operatorname{ad}_{\beta}}\chi(\alpha)
 =
 \frac{\operatorname{ad}_{\beta}}{1-e^{-\operatorname{ad}_{\beta}}}\chi(\alpha).
\]
Expanding the power series
\[
\frac{z}{1-e^{-z}} = 1 + \frac{z}{2} + \frac{z^2}{12} - \frac{z^4}{720} + \cdots
\]
yields~\eqref{eq:beta-bernoulli-fm3}.
\end{proof}

\begin{corollary}[Abelian reduction; \ClaimStatusProvedHere]
\label{cor:fm3-abelian-reduction}
If $\mathfrak g$ is abelian, then all higher brackets vanish,
$\mathfrak P_{\FM_3}(X,\mathfrak g)$ is a strict complex,
and the Maurer--Cartan equation reduces to
\[
 d\alpha = 0,
 \qquad
 d\beta = \chi(\alpha).
\]
Thus the first nontrivial planted-forest correction is genuinely nonabelian.
\end{corollary}

\begin{proof}
If $\mathfrak g$ is abelian, all Lie brackets in
\eqref{eq:l2-fm3}, \eqref{eq:l3-fm3}, and
\eqref{eq:beta-bernoulli-fm3} vanish.
\end{proof}

\section{The first nontrivial Maurer--Cartan obstructions}
\label{sec:fm3-obstructions}

Write a formal degree-one element as
\[
\Theta(t)=t\Theta_1+t^2\Theta_2+t^3\Theta_3+\cdots,
\qquad
\Theta_k=(\alpha_k,\beta_k)\in \mathfrak P_{\FM_3}(X,\mathfrak g)^1.
\]
Expanding the Maurer--Cartan equation
\[
 l_1(\Theta) + \frac12 l_2(\Theta,\Theta) + \frac16 l_3(\Theta,\Theta,\Theta) + \cdots = 0
\]
produces an explicit obstruction tower.

\begin{proposition}[Linear and quadratic obstruction classes; \ClaimStatusProvedHere]
\label{prop:fm3-first-obstructions}
Let $\Theta_1=(\alpha_1,\beta_1)$.
\begin{enumerate}[label=\textup{(\alph*)}]
\item The order-$t$ equations are
  \begin{equation}
  \label{eq:order1-fm3}
  d\alpha_1=0,
  \qquad
  d\beta_1 = \chi(\alpha_1).
  \end{equation}
  Equivalently, the first gluing obstruction is the class
  \begin{equation}
  \label{eq:first-gluing-obstruction}
  o_1(\alpha_1)
  :=
  [\chi(\alpha_1)]
  \in
  H^2\bigl(\mathfrak P_{\FM_3}(X,\mathfrak g),l_1\bigr)
  \cong
  H^1(\mathfrak r_1,d).
  \end{equation}
\item Suppose \eqref{eq:order1-fm3} holds.  Then the order-$t^2$ obstruction is
  the class of
  \begin{equation}
  \label{eq:o2-fm3}
  o_2(\Theta_1)
  :=
  \frac12 l_2(\Theta_1,\Theta_1)
  =
  \left(
  \frac12[\alpha_1,\alpha_1],
  \frac12[\beta_1,\chi(\alpha_1)]
  \right)
  \in
  Z^2\bigl(\mathfrak P_{\FM_3}(X,\mathfrak g),l_1\bigr).
  \end{equation}
  It vanishes in cohomology if and only if there exists
  $\Theta_2=(\alpha_2,\beta_2)$ such that
  \[
  l_1(\Theta_2) + \frac12 l_2(\Theta_1,\Theta_1) = 0.
  \]
\end{enumerate}
\end{proposition}

\begin{proof}
Part~(a) is the order-$t$ part of the Maurer--Cartan equation.  The class
$o_1(\alpha_1)$ is the cohomology class of the degree-two element
$l_1(\alpha_1,0)=(0,\chi(\alpha_1))$.

Part~(b) is the order-$t^2$ part of the Maurer--Cartan equation.  Since
$l_1(\Theta_1)=0$, the $L_\infty$ identities imply that
$l_2(\Theta_1,\Theta_1)$ is $l_1$-closed.  The standard obstruction statement
for filtered Maurer--Cartan equations gives the criterion.
\end{proof}

The class $o_2$ is still ``strictly dg Lie'' in origin: it depends only on $l_1$
and $l_2$.  The first genuinely planted-forest obstruction appears one order
later and is exactly the cubic term forced by the nested boundary.

\begin{theorem}[The first genuinely planted-forest Maurer--Cartan obstruction; \ClaimStatusProvedHere]
\label{thm:fm3-cubic-obstruction}
Assume that $\Theta_1=(\alpha_1,\beta_1)$ and $\Theta_2=(\alpha_2,\beta_2)$
solve the Maurer--Cartan equation modulo $t^3$:
\[
 l_1(\Theta_1)=0,
 \qquad
 l_1(\Theta_2)+\frac12 l_2(\Theta_1,\Theta_1)=0.
\]
Then the obstruction to extending to order $t^3$ is the cohomology class of
\begin{equation}
\label{eq:cubic-obstruction-class}
 o^{\mathrm{pf}}_3(\Theta_1,\Theta_2)
 :=
 l_2(\Theta_1,\Theta_2) + \frac16 l_3(\Theta_1,\Theta_1,\Theta_1)
 \in
 Z^2\bigl(\mathfrak P_{\FM_3}(X,\mathfrak g),l_1\bigr).
\end{equation}
Its divisor component is
\begin{equation}
\label{eq:cubic-obstruction-divisor}
\bigl(o^{\mathrm{pf}}_3(\Theta_1,\Theta_2)\bigr)_{\mathrm{div}}
=
[\alpha_1,\alpha_2],
\end{equation}
and its overlap component is
\begin{equation}
\label{eq:cubic-obstruction-overlap}
\bigl(o^{\mathrm{pf}}_3(\Theta_1,\Theta_2)\bigr)_{\mathrm{ov}}
=
\frac12\Bigl([\beta_1,\chi(\alpha_2)] + [\beta_2,\chi(\alpha_1)]\Bigr)
+ \frac1{12}[\beta_1,[\beta_1,\chi(\alpha_1)]].
\end{equation}
Componentwise, on each nested planted forest $ij\prec 123$ one has
\begin{align}
\label{eq:cubic-obstruction-componentwise}
\bigl(o^{\mathrm{pf}}_3\bigr)_{ij}
&=
\frac12\Bigl([
\beta_{1,ij},\chi_{ij}(\alpha_2)] + [\beta_{2,ij},\chi_{ij}(\alpha_1)]\Bigr)
+ \frac1{12}[\beta_{1,ij},[\beta_{1,ij},\chi_{ij}(\alpha_1)]],
\end{align}
where
\[
\chi_{ij}(\alpha_m)
=
\alpha_{m,ij}|_{S_{ij}} - \alpha_{m,123}|_{S_{ij}}.
\]
The class of~\eqref{eq:cubic-obstruction-class} vanishes if and only if there
exists $\Theta_3$ such that
\[
 l_1(\Theta_3) + l_2(\Theta_1,\Theta_2) + \frac16 l_3(\Theta_1,\Theta_1,\Theta_1)=0.
\]
\end{theorem}

\begin{proof}
This is the order-$t^3$ coefficient of the Maurer--Cartan equation.  Because the
partial solution is valid modulo $t^3$, the general $L_\infty$ obstruction
argument shows that the coefficient
\[
 l_2(\Theta_1,\Theta_2) + \frac16 l_3(\Theta_1,\Theta_1,\Theta_1)
\]
is $l_1$-closed.  Using the explicit formulas
\eqref{eq:l2-fm3} and~\eqref{eq:l3-fm3} proves
\eqref{eq:cubic-obstruction-divisor} and
\eqref{eq:cubic-obstruction-overlap}.  The vanishing criterion is the standard
recursive Maurer--Cartan extension criterion in a filtered complete
$L_\infty$-algebra.
\end{proof}

\begin{remark}[The cubic obstruction is the nested-collision defect]
Equation~\eqref{eq:cubic-obstruction-componentwise} is the first point at which
the geometry of the codimension-two strata is unavoidable.  The first two terms
measure how the second-order divisor correction and the second-order overlap
correction fail to match along $S_{ij}$.  The third term,
\[
\frac1{12}[\beta_{1,ij},[\beta_{1,ij},\chi_{ij}(\alpha_1)]],
\]
is the first genuinely higher correction: it is precisely the binary-in-triple
planted-forest term and has no analogue on $\FM_2(X)$.
\end{remark}

\section{Specialization to the modular deformation algebra}
\label{sec:fm3-specialization-defcyc}

The constructions above are functorial in the coefficient dg Lie algebra
$\mathfrak g$.  Specializing to
\[
\mathfrak g = \Defcyc(\cA)
\]
for a Koszul chiral algebra $\cA$ produces an explicit local small model for the
boundary piece of the cyclic deformation algebra constructed globally in
Theorem~\ref{thm:cyclic-linf-graph}.

\begin{corollary}[The \texorpdfstring{$\FM_3$}{FM3} planted-forest model for modular deformations; \ClaimStatusProvedHere]
\label{cor:fm3-defcyc-specialization}
For every Koszul chiral algebra $\cA$ with cyclic deformation dg Lie algebra
$\Defcyc(\cA)$, the boundary logarithmic deformation algebra of
$\FM_3(X)$ admits a filtered complete small model
\[
\mathfrak P_{\FM_3}(X,\Defcyc(\cA))
=
\operatorname{Cone}(\chi_{\cA})[-1],
\]
where $\chi_{\cA}$ is the residue-incidence morphism obtained by replacing
$\mathfrak g$ with $\Defcyc(\cA)$ in~\eqref{eq:chi-fm3}.  The first higher
operation $l_3$ and the first genuinely planted-forest obstruction are given by
Theorem~\ref{thm:fm3-explicit-brackets} and
Theorem~\ref{thm:fm3-cubic-obstruction}.
\end{corollary}

\begin{proof}
Apply Theorems~\ref{thm:fm3-homotopy-fiber},
\ref{thm:fm3-explicit-brackets}, and~\ref{thm:fm3-cubic-obstruction} with
$\mathfrak g=\Defcyc(\cA)$.
\end{proof}

\section{Synthesis and programme interfaces}
\label{sec:fm3-synthesis-programme}

Nothing in this section is used in the proofs above.  Its purpose is only to
record the exact place of the $\FM_3$ construction in the larger subject that
has been developed throughout the thread.

\begin{enumerate}[label=\textup{(\arabic*)}]
\item \emph{Coderivations to residues.}
  The canonical modular Maurer--Cartan element lives first on the completed bar
  side.  The cyclic deformation package globalizes it.  The present appendix
  isolates the smallest boundary model in which one can see the local residue
  mechanism and the first genuinely higher planted-forest correction.
\item \emph{Worldsheet sewing.}
  The incidence map $\chi$ is the first nontrivial sewing operator: it compares
  residue data on binary-collision divisors with residue data on the triple
  divisor along the nested strata $ij\prec 123$.  The cubic obstruction of
  Theorem~\ref{thm:fm3-cubic-obstruction} is the smallest local shadow of the
  higher-genus modular-operadic sewing constraints.
\item \emph{Open strings and branes.}
  On the boundary side of the string dictionary discussed in the thread, the bar
  construction controls the open sector, while Chan--Paton ordering forces an
  $E_1$ rather than an $E_\infty$ structure.  The present $\FM_3$ model is the
  first local nontrivial test case in which that ordered, noncommutative
  behaviour produces a higher residue correction.
\item \emph{Closed strings.}
  The closed sector is expected to be recovered from a cobar or defect-cyclic
  bulk object.  The overlap equation~\eqref{eq:beta-bernoulli-fm3} identifies
  the exact local gluing problem that any boundary-generated bulk candidate must
  solve along codimension-two nested collisions.
\item \emph{Holography.}
  The cubic planted-forest term is the first place where a local residue defect
  is acted on twice by an overlap parameter.  This is the local prototype for the
  passage from boundary modular data to the nested commutator structures that one
  expects in the line-operator and Yangian-like sectors of twisted holography.
\end{enumerate}

\begin{remark}[What remains after \texorpdfstring{$\FM_3$}{FM3}]
The $\FM_3$ model solves the first nontrivial local problem completely.  The next
step is no longer local contraction; that part is finished.  The next step is
higher-depth gluing:
one must organize the residue dg Lie algebras of all planted forests of $\FM_n(X)$,
transfer them to a global small model, and identify the resulting higher
operations with the universal modular Maurer--Cartan package in the global cyclic
$L_\infty$ algebra.
\end{remark}
